% ==============================================================================
% Fichier : front/02_mode_emploi.tex
% Titre   : Mode d'emploi — TPIN
% ==============================================================================

\chapter*{Mode d'emploi de ce cours}
\addcontentsline{toc}{chapter}{Mode d'emploi}
\markboth{Mode d'emploi}{Mode d'emploi}

\section*{Structure générale}

Ce livre correspond à l'UE \textbf{SEC4031}. Il compte un prologue et trois
parties, soit treize chapitres, pour 60~heures en présentiel réparties en
trente séances de deux heures. Les fiches d'auto-évaluation de NDAP SUKULU
se font hors de ces 60~heures.

\begin{tableaucours}{L{1.4cm}L{5.4cm}L{4.4cm}L{1.6cm}}{%
    \tableauHeader{Partie} &
    \tableauHeader{Chapitres} &
    \tableauHeader{Compétences visées} &
    \tableauHeader{Séances}
}
Prologue & Scène de crime numérique (Ch.~\ref{chap:prologue}) & Appréhender
la réalité du terrain & 1 et 2 \\
I & Philosophie et fondements, histoire, grandes affaires, fondements
théoriques, état de l'art (Ch.~\ref{chap:phi-fond} à~\ref{chap:etat-art}) &
Comprendre l'origine et les bases théoriques de la discipline & 3 à 12 \\
II & Processus investigatif universel, méthodologies et standards, chaîne
de custody, acquisition (Ch.~\ref{chap:PIU} à~\ref{chap:acq}) & Conduire
une investigation structurée et défendable & 13 à 22 \\
III & Cadre normatif global, législation mondiale et régionale, droit
camerounais et africain, déontologie (Ch.~\ref{chap:norm-glob}
à~\ref{chap:deonto}) & Opérer dans la légalité et l'éthique & 23 à 30 \\
\end{tableaucours}

\begin{boxinfo}
\textbf{La suite~: l'UE SEC4052.} \emph{Investigation Numérique Avancée}
prolonge ce cours~: forensique opérationnelle (système, réseau, mobile,
nuage, objets connectés, sources ouvertes, rapport, cas intégrateurs),
anti-forensique et défis émergents, puis fondements cryptographiques de la
preuve. Elle suppose acquis les gestes et les principes de SEC4031.
\end{boxinfo}

\section*{Les environnements visuels}

\begin{boxdef}
\textbf{Définition}. Introduit un terme technique avec sa définition formelle.
Tous les termes définis sont indexés en fin d'ouvrage.
\end{boxdef}

\begin{boxobjectifs}
Objectifs du chapitre. Liste les compétences mesurables attendues à l'issue
du chapitre. Ces objectifs correspondent aux critères d'évaluation.
\end{boxobjectifs}

\begin{boxCRO}
Analyse CRO. Évalue systématiquement une technique ou un protocole selon
les trois dimensions du Trilemme : Confidentialité~($\mathcal{C}$),
Fiabilité~($\mathcal{R}$), Opposabilité~($\mathcal{O}$).
\end{boxCRO}

\begin{boxscenario}{Titre du scénario}
Scénario opérationnel. Présente un cas réel ou fictif ancré dans la réalité
camerounaise ou africaine. Le cas MBONGO (fraude financière 12M~FCFA) sert
de fil rouge du prologue à la Partie~II, puis dans le cours SEC4052.
\end{boxscenario}

\begin{boxPNC}
Application OPJ/DGSN. Traduit le concept académique en procédure opérationnelle
destinée aux Officiers de Police Judiciaire.
\end{boxPNC}

\begin{boxcustody}
Chaîne de Custody. Détaille les étapes et formulaires nécessaires à la
préservation de la valeur probante d'une preuve numérique.
\end{boxcustody}

\begin{boxartefact}{Nom de l'artefact}
Artefact forensique. Décrit un artefact numérique spécifique : sa localisation
sur le système, la valeur qu'il porte pour l'investigation, et sa fragilité.
\end{boxartefact}

\begin{boxjuris}
Jurisprudence \& Droit applicable. Référence les textes de loi, articles et
jurisprudences pertinents, avec accent sur le droit camerounais et africain.
\end{boxjuris}

\begin{boxalert}
Attention. signale un risque procédural, légal ou technique pouvant
invalider une investigation.
\end{boxalert}

\begin{boxresume}
À retenir. Synthèse des points essentiels du chapitre.
\end{boxresume}

\begin{boxexo}
Exercice. Travaux pratiques gradués. Chaque exercice porte un indicateur
de niveau : \niveaubase\ fondamental, \niveaumoyen\ intermédiaire,
\niveauexpert\ expert.
\end{boxexo}

\section*{Conventions typographiques}

\begin{itemize}
    \item Les \termecle{termes clés} apparaissent en \textbf{gras bleu marine}
          et sont indexés automatiquement.
    \item Les \sigle{ACRONYMES} utilisent des petites capitales.
    \item Les références aux normes utilisent la commande \norme{ISO/IEC~27037}.
    \item Les références aux lois camerounaises utilisent \loicam{2010/012}.
    \item Les hachages cryptographiques sont formatés ainsi :
          \hash{SHA-256}{\texttt{3a4b5c6d...}}.
    \item L'impact CRO inline est noté : \cro{+0.8}{-0.3}{+0.6}.
    \item L'ordre de volatilité est indiqué : \volatilite{1}{RAM}.
\end{itemize}

\section*{Matériel d'accompagnement}

Ce cours est complété par :
\begin{itemize}
    \item \textbf{Images de pratique forensique}, images E01 préparées pour
          les TP Autopsy.
    \item \textbf{Captures réseau (.pcap)}, traces Wireshark pour les TP
          de forensique réseau.
    \item \textbf{Machines virtuelles}, VM Kali Linux et VM Windows avec
          artefacts préinstallés (clé USB fournie en TP).
    \item \textbf{Formulaires officiels}, formulaire de chaîne de custody
          (Annexe~C), adapté du formulaire type DGSN.
\end{itemize}
